%==============================================================
%  lecons/ensembles/construction_N.tex — Construction de N
%==============================================================

\lecontitle{Construction de \texorpdfstring{$\mathbb{N}$}{N}}{Théorie des ensembles — Axiomes de Peano}

%=============================================================
%  § 1 — MOTIVATION ET AXIOMES DE PEANO
%=============================================================
\secnum{navyblue}{1}{Axiomes de Peano}

\begin{tcolorbox}[thmbox]
  \textbf{Axiomes de Peano (1889).}\enspace%
  \index{Peano (axiomes de)}\index{entiers naturels!axiomes de Peano}
  Il existe un ensemble $\mathbb{N}$ muni d'un élément distingué $0 \in \mathbb{N}$
  et d'une application $S : \mathbb{N} \to \mathbb{N}$ (le \emph{successeur}) vérifiant :
  \begin{enumerate}
    \item[\textbf{P1.}] $0$ n'est le successeur d'aucun entier : $0 \notin \mathrm{Im}(S)$.
    \item[\textbf{P2.}] $S$ est injective : $S(n) = S(m) \Rightarrow n = m$.
    \item[\textbf{P3.}] \textbf{(Récurrence)} Si $A \subseteq \mathbb{N}$ vérifie
      $0 \in A$ et $\forall n \in A,\; S(n) \in A$, alors $A = \mathbb{N}$.
  \end{enumerate}
  On pose $1 = S(0)$, $2 = S(1)$, $3 = S(2)$, etc.
\end{tcolorbox}

\rubriq{Unicité à isomorphisme près}
Si $(\mathbb{N}, 0, S)$ et $(\mathbb{N}', 0', S')$ satisfont tous deux les axiomes
de Peano, il existe un unique isomorphisme $\varphi : \mathbb{N} \to \mathbb{N}'$
tel que $\varphi(0) = 0'$ et $\varphi \circ S = S' \circ \varphi$.
La preuve utilise P3 deux fois : d'abord pour construire $\varphi$ par récurrence,
ensuite pour montrer que c'est une bijection.

\decsep{}

%=============================================================
%  § 2 — ADDITION ET MULTIPLICATION
%=============================================================
\secnum{navyblue}{2}{Définition de l'addition et de la multiplication}

\begin{tcolorbox}[thmbox]
  \textbf{Définition par récurrence.}\enspace%
  \index{entiers naturels!addition}\index{entiers naturels!multiplication}
  On définit $+$ et $\times$ sur $\mathbb{N}$ par :
  \[
    n + 0 = n, \qquad n + S(m) = S(n + m),
  \]
  \[
    n \times 0 = 0, \qquad n \times S(m) = (n \times m) + n.
  \]
  Ces définitions sont légitimes par le \emph{théorème de définition par récurrence} :
  pour toute application $f : \mathbb{N} \to \mathbb{N}$ et tout $a \in \mathbb{N}$,
  il existe une unique suite $(u_n)$ telle que $u_0 = a$ et $u_{S(n)} = f(u_n)$.
\end{tcolorbox}

\begin{mdframed}[style=proofbar]
  On vérifie par récurrence les propriétés algébriques usuelles :
  commutativité, associativité de $+$ et $\times$, distributivité.
  Par exemple, $n + m = m + n$ se montre par double récurrence sur $n$ et $m$.

  \textit{Ordre.} On pose $n \leq m$ si $\exists k \in \mathbb{N},\; m = n + k$.
  C'est un ordre total (P3 garantit la comparabilité) et compatible avec $+$ et $\times$.
  $\blacksquare$
\end{mdframed}

\decsep{}

%=============================================================
%  § 3 — MODÈLE DE VON NEUMANN
%=============================================================
\secnum{navyblue}{3}{Modèle ensembliste — construction de von Neumann}

\begin{tcolorbox}[thmbox]
  \textbf{Construction de von Neumann (dans ZFC).}\enspace%
  \index{von Neumann (construction de)}\index{ZFC}
  On construit $\mathbb{N}$ comme ensemble au sein de la théorie des ensembles
  de Zermelo-Fraenkel avec l'axiome du choix (ZFC) :
  \[
    0 := \emptyset, \quad S(n) := n \cup \{n\}, \quad
    \mathbb{N} := \{0,\, S(0),\, S(S(0)),\, \ldots\}.
  \]
  Explicitement : $0=\emptyset$, $1=\{\emptyset\}$, $2=\{\emptyset,\{\emptyset\}\}$,
  $3=\{\emptyset,\{\emptyset\},\{\emptyset,\{\emptyset\}\}\}$, etc.
  L'existence de $\mathbb{N}$ est garantie par l'\emph{axiome de l'infini}.
\end{tcolorbox}

\rubriq{Signification}
Chaque entier $n$ est identifié à l'ensemble $\{0,1,\ldots,n-1\}$ de ses
prédécesseurs. Ainsi $n \in m \Leftrightarrow n < m$, ce qui donne une
interprétation purement ensembliste de l'ordre.

\decsep{}

%=============================================================
%  § 4 — EXEMPLES, PIÈGES ET EXERCICES
%=============================================================
\secnum{exgreen}{4}{Propriétés, pièges et exercices}

\rubriq{Propriétés fondamentales}

\begin{mdframed}[style=exbar]
  \textbf{1. Principe de bonne fondation.}\enspace%
  \index{bonne fondation}
  Tout sous-ensemble non vide $A \subseteq \mathbb{N}$ admet un plus petit élément.
  \textit{Preuve} : on utilise le lemme de l'ordre discret
  \[
    \forall a,n\in\mathbb{N},\qquad n < a \;\Longrightarrow\; S(n) \leq a,
  \]
  qui exprime qu'aucun entier ne se glisse strictement entre $n$ et $S(n)$
  (il se démontre par récurrence sur $a$ à partir de la définition de $\leq$).
  Soit $A$ non vide et supposons par l'absurde qu'il n'a pas de plus petit
  élément. Posons
  \[
    B = \{n\in\mathbb{N} \mid n \leq a \text{ pour tout } a\in A\}.
  \]
  \emph{Initialisation :} $0 \leq a$ pour tout $a$, donc $0\in B$.
  \emph{Hérédité :} soit $n\in B$, c'est-à-dire $n\leq a$ pour tout $a\in A$.
  Si l'on avait $n\in A$, alors $n$ serait un minorant de $A$ appartenant à $A$,
  donc le plus petit élément de $A$, ce qui est exclu ; ainsi $n\notin A$,
  et pour tout $a\in A$ l'inégalité $n\leq a$ est stricte : $n<a$.
  Le lemme donne alors $S(n)\leq a$ pour tout $a\in A$, donc $S(n)\in B$.
  Par P3, $B = \mathbb{N}$. Comme $A$ est non vide, prenons $a_0\in A$ ;
  alors $S(a_0)\in B$ entraîne $S(a_0)\leq a_0$, ce qui est absurde
  (on a toujours $a_0 < S(a_0)$). $A$ admet donc un plus petit élément.

  \smallskip\noindent
  \textbf{2. Récurrence forte.}\enspace%
  Si $P(n)$ est vraie dès que $P(k)$ est vraie pour tout $k < n$,
  alors $P(n)$ est vraie pour tout $n$. Équivalent à P3 via la bonne fondation.

  \smallskip\noindent
  \textbf{3. Cardinalité.}\enspace%
  $\mathbb{N}$ est infini dénombrable : $|\mathbb{N}| = \aleph_0$.
  Tout sous-ensemble infini de $\mathbb{N}$ est dénombrable.
\end{mdframed}

\vspace{6pt}
\rubriq{Pièges}

\begin{mdframed}[style=cexbar]
  \textbf{$0 \in \mathbb{N}$ ou $0 \notin \mathbb{N}$ ?}\enspace%
  Par convention (ISO 80000-2, Peano), $0 \in \mathbb{N}$. Certains auteurs
  (tradition française classique) posent $\mathbb{N} = \{1,2,3,\ldots\}$.
  Vérifier la convention avant tout raisonnement.

  \smallskip\noindent
  \textbf{Peano au premier ordre : incomplétude et modèles non standard.}\enspace%
  Au premier ordre, l'axiome de récurrence (P3) devient un \emph{schéma} :
  il ne porte que sur les parties de $\mathbb{N}$ définissables par une formule,
  et non sur toutes les parties (catégoricité perdue par rapport à la version
  du second ordre). Deux conséquences classiques en découlent.
  D'une part, par le théorème de \emph{compacité} (ou Löwenheim--Skolem),
  l'arithmétique de Peano admet, à côté du modèle standard $\omega$, des
  \emph{modèles non standard} (contenant des entiers «\,infinis\,») satisfaisant
  exactement les mêmes énoncés du premier ordre.
  D'autre part, par le théorème d'\emph{incomplétude} de Gödel, si PA est
  cohérente il existe des énoncés arithmétiques vrais (dans $\omega$) mais
  non démontrables dans PA.
\end{mdframed}

\vspace{6pt}
\exolabel

\begin{mdframed}[style=exobar]
  \begin{enumerate}
    \item Montrer par récurrence que $\sum_{k=0}^{n} k = n(n+1)/2$.

    \item En utilisant les axiomes de Peano, montrer que la soustraction
          $m - n$ (pour $n \leq m$) est bien définie et unique.

    \item Montrer que $\mathbb{N}$ est le seul modèle à isomorphisme près
          des axiomes de Peano (unicité du modèle standard).

    \item (Division euclidienne) Montrer que pour tous $n,d \in \mathbb{N}$
          avec $d > 0$, il existe des uniques $q, r \in \mathbb{N}$ tels que
          $n = qd + r$ et $0 \leq r < d$.

    \item Montrer que dans la construction de von Neumann, $n \in m
          \Leftrightarrow n \subsetneq m \Leftrightarrow n < m$.
  \end{enumerate}
\end{mdframed}

\leconfooter{G.\ Peano (1889) — J.\ von Neumann (1923) — E.\ Zermelo (1908) \& A.\ Fraenkel (1922)}
